\begin{derivation}[外场做功与近场复耦合]
采用全文的正频率复幅约定，把真实外场写成
\[
E_{{\rm ext},z}(\vb{r},t)
=\int_0^\infty\frac{\dd\omega}{2\pi}
\left[
E_{{\rm ext},\omega,z}(\vb{r})\ee^{-\ii\omega t}
+E_{{\rm ext},\omega,z}^*(\vb{r})\ee^{+\ii\omega t}
\right].
\]
沿匀速轨迹 $\vb{r}_e(t)=(\vb{R}_0,v_0t)$ 代入正文\eqnrefs{eq:pinem_time_domain_work_dp68}。对固定正频率 $\omega$，定义复功幅
\begin{align*}
\mathcal W_\omega
&\equiv
q_{\mathrm e}\int_{-\infty}^{\infty}\dd t\,
 v_0 E_{{\rm ext},\omega,z}(\vb{R}_0,v_0t)
 \ee^{-\ii\omega t}\\
&=q_{\mathrm e}\int_{-\infty}^{\infty}\dd z\,
 E_{{\rm ext},\omega,z}(\vb{R}_0,z)
 \ee^{-\ii\omega z/v_0},
\end{align*}
其中第二行只用了 $z=v_0t$ 与 $\dd z=v_0\dd t$。指数因子表明电子只抽取外场纵向 Fourier 谱在
\[
k_z=\frac{\omega}{v_0}
\]
这一同步线上对应的分量。

单个边带步长的能量是 $\hbar\omega$，因此把相干功幅除以这一能量定义无量纲生成元
\[
\beta(\omega)\equiv\frac{\mathcal W_\omega}{\hbar\omega},
\]
立刻得到正文\eqnrefs{eq:pinem_beta_green_unified_dp47}。这个除法不是经验归一化：在能量格演化算符中，$\beta$ 必须是乘在无量纲平移生成元前的无量纲系数。

若外场由外源电流产生，正文采用
\[
\vb{E}_{{\rm ext},\omega}(\vb{r})
=\ii\mu_0\omega
\int\dd^3r'\,
\mathbf G^R(\vb{r},\vb{r}';\omega)
\cdot\widetilde{\vb{J}}_{\rm ext}(\vb{r}',\omega).
\]
把它代回 $\beta$ 并约去 $\omega$，得到
\begin{align*}
\beta(\omega)
={}&\frac{\ii\mu_0q_{\mathrm e}}{\hbar}
\int\dd z\dd^3r'\,
 \ee^{-\ii\omega z/v_0}
\hat{\vb{z}}\cdot
\mathbf G^R((\vb{R}_0,z),\vb{r}';\omega)
\cdot\widetilde{\vb{J}}_{\rm ext}(\vb{r}',\omega).
\end{align*}
因此 PINEM 耦合与 EELS/CL 使用的是同一个推迟 Green 算符；PINEM 由外源产生相干场并沿电子轨迹读取线性振幅，EELS 则对电子自身诱导场取耗散二次型。
\end{derivation}
